\section{Definição do Sistema}

\subsection{Objetivo do Sistema}

O sistema tem como objetivo auxiliar o usuário na análise e seleção
de carteiras de investimentos, considerando a relação entre risco e
retorno dos ativos selecionados.

A partir dos dados dos ativos, o sistema deverá avaliar diferentes
composições de carteira e identificar aquelas que apresentam uma
relação eficiente entre risco e retorno segundo o modelo de
Markowitz.

\subsection{Escopo}

O sistema será destinado à análise de carteiras compostas por
múltiplos ativos. A análise deverá considerar os dados dos ativos
selecionados para calcular o retorno esperado, o risco e a relação
entre os ativos.

O sistema deverá permitir a análise de até 30 ativos
simultaneamente em uma mesma carteira.

Como resultado, o sistema deverá apresentar informações sobre a
composição das carteiras analisadas, seus respectivos riscos e
retornos, além da representação da fronteira eficiente.

\subsection{Perfil do Usuário}

O sistema será utilizado por investidores ou analistas interessados
em analisar diferentes combinações de ativos e avaliar alternativas
de composição de carteiras.

O usuário será responsável por selecionar os ativos e fornecer ou
disponibilizar os dados necessários para a realização da análise.

\subsection{Necessidades do Sistema}

Para atender ao objetivo proposto, o sistema deverá possibilitar:

\begin{itemize}
\item a análise conjunta de diferentes ativos;


\item a avaliação da relação entre risco e retorno;

\item a comparação entre diferentes composições de carteira;

\item a identificação de carteiras eficientes;

\item a consulta dos pesos atribuídos aos ativos;

\item a apresentação dos resultados por meio de tabelas e
gráficos.


\end{itemize}

\subsection{Conceitos do Domínio}

O sistema está inserido no domínio de análise e seleção de
carteiras de investimentos. Os principais conceitos considerados
são:

\begin{itemize}
\item \textbf{Ativo:} instrumento financeiro que pode fazer
parte de uma carteira;


\item \textbf{Carteira:} conjunto de ativos organizados de
acordo com determinada distribuição de pesos;

\item \textbf{Peso:} participação de cada ativo na composição
da carteira;

\item \textbf{Retorno esperado:} estimativa do retorno de um
ativo ou de uma carteira;

\item \textbf{Risco:} medida da variabilidade dos retornos da
carteira;

\item \textbf{Covariância:} medida que representa a relação
entre os retornos de dois ativos;

\item \textbf{Fronteira eficiente:} conjunto de carteiras que
apresentam as melhores combinações entre risco e retorno,
considerando o critério de eficiência adotado pelo modelo de
Markowitz.


\end{itemize}

\subsection{Regras Gerais}

A análise das carteiras será baseada no modelo de seleção de
carteiras de Markowitz.

Considerando a restrição de não permitir vendas a descoberto, os
pesos dos ativos deverão satisfazer:

\begin{equation}
X_i \geq 0
\end{equation}

Além disso, a soma dos pesos dos ativos que compõem uma carteira
deverá ser igual a 1:

\begin{equation}
\sum_{i=1}^{n} X_i = 1
\end{equation}

A análise deverá considerar o retorno esperado dos ativos e a
covariância entre eles para determinar o risco das diferentes
composições de carteira.

\subsection{Limitações e Desempenho}

O sistema deverá considerar como limite operacional a análise
simultânea de até 30 ativos em uma mesma carteira.

Para esse limite máximo, o sistema deverá realizar os cálculos
necessários e apresentar a representação gráfica dos resultados em
até 2 segundos.
